\chapter{Langues peu dotées}

Les langues peu dotées sont au centre de la recherche menée dans ce mémoire. Ces langues (corse,
occitan et alsacien) constituent des cas d’étude particulièrement pertinents, car elles combinent
plusieurs caractéristiques~: faible disponibilité de ressources numériques, variations dialectales
importantes, normalisation orthographique parfois incomplète et faible présence dans les grands
jeux de données multilingues utilisés en TAL. Elles représentent ainsi un terrain d’expérimentation
idéal pour les méthodes adaptées aux environnements peu dotés. Elles posent
diverses problématiques qui, si résolues, permettraient de faire avancer la recherche sur le TAL.

Hélas, il est important de noter que les travaux de recherche en TAL, dans leur majorité, se sont concentrés
sur des langues, généralement bien décrites et étant majoritaires ou parlées par un grand nombre de locuteurs \citep{delgratta_2014, bernhard_2018}.

Cependant de plus en plus de travaux de recherche universitaire de grande qualité sont menés
pour les langues peu dotées. Ce qui pourrait être lié à une prise de conscience majeure et partagée dans le domaine du TAL
sur l’importance de la représentation numérique pour toutes les langues \citep{bernhard_2018}.

Les langues régionales étudiées dans ce mémoire ne présentent que peu de ressources numériques et sont par définition sous-dotées. Les ressources numériques
existantes sont créées uniquement par les communautés utilisant encore ces langues. Or, celles-ci sont en déclin et de moins
en moins de personnes décident d'apprendre et d'utiliser les langues régionales. Dans le cas de l'alsacien, par exemple, après la seconde
guerre mondiale, en 1946, 90,8~\% de la population d'Alsace se déclarait dialectophone. De nos jours, ce nombre est tombé
(à compter de 2022) à 46~\% d'après l'\href{https://www.olcalsace.org/}{OLCA}\footnote{
    Nombres de l'OLCA~: \href{https://www.olcalsace.org/fr/observer-et-veiller/le-dialecte-en-chiffres}{le dialecte en chiffres}}.
Cette baisse peut être expliquée par une utilisation du français comme
langue principale partout au sein du pays et par une perte des langues régionales au fil du temps au profit d'une langue
française généralisée.

\section{Définition}

Une langue peu dotée peut être définie comme une langue disposant d’un nombre limité de ressources
linguistiques numériques exploitables pour des tâches de TAL. Cette sous-dotation peut concerner l’absence
de corpus annotés, de lexiques structurés, de modèles pré-entraînés ou de normes orthographiques stabilisées.
Elle ne recouvre pas nécessairement la notion de langue minoritaire. Certaines langues peuvent être
majoritaires dans un pays mais restent sous-dotées \citep{gerken_2022}.
Il est à noter que certaines langues, pourtant minoritaires ou menacées, ne seront pas sous-dotées. Le catalan, par exemple,
rentre dans cette catégorie de langues minoritaires qui ne sont pourtant pas sous-dotées \citep{bernhard_2018}.

Il est aussi important de remarquer qu'une langue peu dotée ou minoritaire peut être parlée par beaucoup,
mais posséder peu voire aucune production écrite. C'est cette absence de production écrite qui peut faire que la langue est
peu dotée. Elle devient donc complexe à étudier, particulièrement dans le cadre de la recherche. Mais celle-ci
reste intéressante à étudier de par ses implications pour le TAL. Une langue peu dotée pose plusieurs
problèmes importants qui, une fois résolus, permettront un meilleur traitement de celle-ci par l'informatique.
Mais aussi un meilleur traitement des langues plus communes de par les recherches menées sur des langues
possédant moins de données informatisées.

\section{Problèmes} \label{problemes_peudot}

Les langues peu dotées posent les problèmes suivants~:
\begin{enumerate}
\item Absence partielle ou totale de documents écrits informatisés rendant l'élaboration de corpus parallèles
complexe voire impossible.
\item Recherche dans les langues sous-dotées réalisée par des groupes isolés ou dispersés
\citep{bernhard_2018}. Ces travaux, non centralisés, complexifient la recherche dans le domaine.
\item Variations dans la manière d'écrire une langue, dues soit à une évolution des normes orthographiques dans le temps,
soit à l’absence de ces normes pour les langues essentiellement orales \citep{kurimo_2017} \citep{bernhard_2018}
\item Une bonne performance est compliquée à atteindre sur des systèmes de traduction par réseau
neuronal pour des langues peu dotées en raison de leur nature et l'absence de paires de phrases bilingues de
bonne qualité en grande quantité. \citep{steingrimsson_2023}
\item Manque d'intérêt de la part de la communauté scientifique pour différentes raisons~: «~Manque de
financement, de ressources humaines, de technologie appropriée, de descriptions linguistiques complètes et
précises, de reconnaissance académique~» \citep{bernhard_2018}
\item Un pourcentage important de segments bruités ou mal formés (erreurs typographiques, mélanges de langues,
traductions partielles, segments tronqués) peut être observé dans les ressources disponibles. Ces segments
dégradent la qualité des modèles d’alignement et compliquent l’extraction automatique de correspondances fiables.
dans le peu d'écrits disponibles. Ceci pourrait affecter la performance de logiciels d'alignement. \citep{steingrimsson_2023}
\end{enumerate}

Ces problèmes sont particuliers, car ils touchent la vaste majorité des langues qui sont considérées comme «~peu dotées~»
parmi toutes les langues existantes dans le monde de nos jours \citep{sharoff_2016}.
Cela implique que les langues peu dotées, parlées par environ 3 milliards de personnes, sont peu voire pas étudiées.
Ces problèmes commencent avec le temps à avoir diverses solutions intéressantes, passant souvent par la
recherche en traduction automatique ou en TAL.

\begin{figure}[H]
    \centering
    \includegraphics[width=0.65\textwidth]{joshi_et_al_labels.png}
    \\
    \smallskip
    \small{La taille des cercles pleins colorés réprésentent le nombre de langues présentes dans un groupe.\\}
    \small{Les bulles, en pointillées, délimitent les 6 catégories.}
    \smallskip
    \caption{Nuage de points des labels établis par \citet{joshi_2020}.}
    \label{joshi_et_al_labels_fig}
\end{figure}

D'après le graphique de la figure \ref{joshi_et_al_labels_fig}, on peut remarquer que la vaste majorité des langues existantes se partagent
entre la catégorie 0 et 1. D'après \citet{joshi_2020} les langues rentrent dans ces
6 catégories\footnote{"0 The Left-Behinds, 1 The Scraping-Bys, 2 The Hopefuls, 3 The Rising Stars, 4 The Underdogs, 5 The Winners" nom original
des catégories, notre traduction apparaît dans l'énumération}~:
\begin{enumerate}
    \setcounter{enumi}{-1}
    \item Les abandonnées~: langues ne possédant quasiment aucune ressource, même non annotées, une absence de corpus brut empêche un
    travail correct sur ses langues
    \item Celles qui s'en sortent~: langues ne possédant que peu de ressources non annotées et quasiment aucune ressource annotée,
    des travaux en TAL pour la conservation de ces langues existent
    \item Les optimistes~: possèdent un petit set de données non annotées et annotées et des communautés plus grandes
    \item Les étoiles montantes~: pas assez de données annotées mais une présence en ligne forte à travers les communautés de locuteurs
    \item Les outsiders~: capables de s'en sortir avec des communautés TAL dédiées qui s'en occupent
    \item Les vainqueurs~: présence en ligne dominante, possédant le plus de ressources
\end{enumerate}
Les langues régionales que nous étudions se situent dans des niveaux entre 0 et 1.
Le français, lui, se situe au niveau 5 \citep{stosic_2024}. Cette taxonomie, se basant sur le nombre de ressources langagières
existantes, nous permet de déterminer qu'il est bien plus dur de trouver des ressources correctes permettant
d'élaborer un corpus parallèle pour les langues régionales qu'un corpus comparable.

\section{Solutions explorées}

Pour commencer, nous pouvons remarquer une augmentation de la recherche dans le sujet des langues peu
dotées importantes, qui permet d'observer plusieurs nouvelles solutions explorées par différents groupes de chercheurs.

\newpage

\subsection{Élaboration d'un corpus parallèle}

On peut citer le travail sur l'élaboration d'un corpus parallèle pour les langues
régionales de France avec le corpus parallèle ParCoLab \citep{stosic_2024}. Ce projet vise à créer
un corpus parallèle pour des langues peu dotées comme certaines langues régionales de France
(occitan, corse, alsacien ou poitevin saintongeais).

\subsection{Ressources collaboratives}

\citet{miletic_2018} proposent par exemple la composition de larges corpus annotés pour le serbe en adaptant des ressources de
langues proches, en utilisant des ressources lexicales collaboratives, comme Wiktionary, ou en pré-annotant
le corpus automatiquement. Il est à noter qu'en 2018 d'autres outils pour le traitement du serbe existaient
déjà, mais que la plupart n'étaient pas disponibles librement.

\citet{millour_2018}, quant à elles, proposent une méthodologie par l'annotation en parties de discours basée sur le crowdsourcing dans le cas de l'alsacien et du créole guadeloupéen.
Les chercheuses ont mis en place 2 plateformes permettant le crowdsourcing. Cette méthode consiste à demander à un grand nombre de
personnes de contribuer à un projet, en général via internet. Une des plateformes possède des corpus pré-annotés
présentés aux participants avant d'entraîner des étiqueteurs et l'autre permet une collecte de corpus bruts.

\subsection{Intelligence artificielle pour les langues peu dotées}

Le comité \textit{ad hoc} sur l’intelligence artificielle (\gls{CAHAI}) du Conseil de l’Europe
«~a constaté que l'IA peut aider les locuteurs de langues minoritaires à participer
plus activement à la vie publique, aux débats, aux délibérations ou aux décisions, mais aussi mieux
faire connaître ces langues.~» \citep{gerken_2022} De plus, d'après \citet{gerken_2022}, le CAHAI aurait remarqué que
cela permet de faire mieux connaître ces langues, amenant plus tard, à une meilleure représentation de
ces langues en ligne. L'IA et les réseaux neuronaux étant en plein essor,
il est intéressant de considérer cette possibilité dans le cas de la création de ressources pour des
langues peu dotées.

\newpage

\section{Bonnes nouvelles}

À l'ère d'Internet, il est important de noter que de plus en plus d'entre nous utilisons le web. Les
utilisateurs, parlant des langues minoritaires ou peu dotées, participent parfois à des efforts
d'écriture communautaire et ajoutent leur langue sur des sites où celle-ci devient alors visible à tous.
Cela est très intéressant, car la diffusion d’une langue régionale ou minoritaire sur les plateformes en
ligne contribue à renforcer sa visibilité. C’est notamment le cas du catalan, qui figure parmi les
vingt langues majoritairement traduites sur Wikipédia \citep{gerken_2022}.

On peut aussi remarquer que l'Union européenne essaye de rendre l'emploi des langues minoritaires ou
régionales plus fréquentes dans des domaines importants d'après \citet{gerken_2022}~: l'administration,
l'éducation, la justice ou les médias. Cet effort est exacerbé par la charte européenne des langues
régionales ou minoritaires «~ratifiée par 25 États, et signée (mais non ratifiée) par neuf autres~»
\citep{gerken_2022}. La France signe cette charte en 1999 mais ne l'a toujours pas ratifiée. Celle-ci n'étant
pas ratifiée, la France n'a mis en place aucun engagement.

\vspace{5mm}

Ainsi, malgré un intérêt croissant pour les langues peu dotées, la rareté et l’hétérogénéité des ressources disponibles
rendent nécessaire l’exploration de stratégies alternatives à la constitution directe de corpus parallèles.
Dans cette perspective, les corpus comparables apparaissent comme une piste prometteuse, ce qui conduit à examiner
plus précisément leur définition, leurs propriétés et leurs usages en TAL.
